\section{Deployment with Docker}
\label{sec:deployment}

\subsection{Containerization Strategy}

Containerization ensures the application runs identically across development, testing, and production 
environments. Docker provides:

\begin{enumerate}
  \item \textbf{Consistency}: Identical runtime environment across machines.
  
  \item \textbf{Isolation}: Application dependencies don't interfere with host or other containers.
  
  \item \textbf{Portability}: Images can be deployed on any system with Docker.
  
  \item \textbf{Scalability}: Orchestrators (Kubernetes, Swarm) manage multiple container instances.
  
  \item \textbf{Reproducibility}: Applications are versioned and immutable.
\end{enumerate}

\subsection{Dockerfile Design}

The Dockerfile follows best practices:

\subsubsection{Multi-Stage Builds (Future Enhancement)}

The current Dockerfile uses a single stage. For production, multi-stage builds can reduce image size:

\begin{verbatim}
  FROM python:3.10-slim AS builder
  # Install dependencies, build wheels
  
  FROM python:3.10-slim
  # Copy only wheels, reduce final image size
\end{verbatim}

This is deferred to production deployment.

\subsubsection{Layer Caching Optimization}

Layers are ordered to maximize cache efficiency:

\begin{enumerate}
  \item Install system packages (rarely changes).
  \item Copy \code{requirements.txt} (changes less frequently).
  \item Install Python dependencies (influenced by requirements.txt).
  \item Copy application code (changes frequently).
\end{enumerate}

This ordering ensures that code changes don't invalidate the expensive dependency installation layer.

\subsubsection{Image Size and Security}

The Dockerfile uses \code{python:3.10-slim}, which is $\approx 150$ MB (vs.~$\approx 900$ MB for \code{python:3.10}). 
Minimal base images reduce:

\begin{itemize}
  \item Deployment time (less data to transfer).
  \item Storage requirements.
  \item Attack surface (fewer system tools).
\end{itemize}

\subsection{Health Checks}

The Dockerfile includes a HEALTHCHECK directive:

\begin{verbatim}
  HEALTHCHECK --interval=30s --timeout=10s 
    --start-period=5s --retries=3 \
    CMD curl -fsS http://localhost:8000/health || exit 1
\end{verbatim}

\textbf{Parameters:}
\begin{itemize}
  \item \code{--interval=30s}: Check health every 30 seconds.
  \item \code{--timeout=10s}: Consider check failed if no response in 10 seconds.
  \item \code{--start-period=5s}: Allow 5 seconds for application startup before checking.
  \item \code{--retries=3}: Mark container unhealthy after 3 consecutive failed checks.
\end{itemize}

\textbf{Health Check Logic:}
The check calls \code{curl} to the \code{/health} endpoint. Exit code 0 (success) marks the container healthy; 
any other exit code marks it unhealthy.

Container orchestrators use health status to:
\begin{itemize}
  \item Restart unhealthy containers.
  \item Route traffic away from degraded instances.
  \item Trigger alerts to operators.
\end{itemize}

\subsection{Docker Compose Orchestration}

Docker Compose simplifies multi-container deployments and local testing:

\subsubsection{Service Definition}

\begin{verbatim}
  version: '3.8'
  
  services:
    api:
      build: .                           # Build from Dockerfile
      ports:
        - "8000:8000"                    # Map port 8000
      volumes:
        - ./models:/app/models           # Mount model artifact
      environment:
        - MODEL_PATH=/app/models/svd_model.pkl  # Config
      restart: unless-stopped            # Auto-restart policy
\end{verbatim}

\subsubsection{Volume Management}

The \code{./models:/app/models} volume mount enables:

\begin{itemize}
  \item \textbf{Shared state}: The model artifact persists across container restarts.
  
  \item \textbf{Model updates}: New model artifacts can be deployed by updating the local file 
  (no container rebuild).
  
  \item \textbf{Development workflow}: Local changes to the model are immediately reflected in the 
  running container.
\end{itemize}

\subsubsection{Restart Policy}

The \code{restart: unless-stopped} policy ensures:

\begin{itemize}
  \item \textbf{Resilience}: If the application crashes, Docker automatically restarts it.
  
  \item \textbf{Control}: Explicit \code{docker compose down} stops all services (not restarted).
  
  \item \textbf{Development-friendly}: Unlike \code{restart: always}, this allows developers to stop 
  services without fighting automatic restarts.
\end{itemize}

\subsection{Environment Variables}

Configuration is managed through environment variables, enabling deployment flexibility:

\begin{itemize}
  \item \code{MODEL\_PATH}: Path to the pickled model (default: \code{models/svd\_model.pkl}).
  \item \code{MODEL\_VERSION}: Version string returned in responses (default: \code{1.0.0}).
  \item \code{HOST}: Bind address for the API server (default: \code{0.0.0.0}).
  \item \code{PORT}: Bind port (default: \code{8000}).
  \item \code{DEBUG}: Enable debug mode (default: \code{false}).
\end{itemize}

This design enables:
\begin{itemize}
  \item Different configurations for dev/staging/production without code changes.
  \item Secrets management outside of code/images (credentials can be injected at runtime).
  \item A/B testing by deploying multiple versions with different \code{MODEL\_VERSION} values.
\end{itemize}

\subsection{Deployment Workflow}

A typical deployment follows this workflow:

\begin{enumerate}
  \item \textbf{Build}: \code{docker compose build} creates the image.
  
  \item \textbf{Start}: \code{docker compose up -d} starts the container in background.
  
  \item \textbf{Verify}: \code{curl http://localhost:8000/health} checks health.
  
  \item \textbf{Test}: Run integration tests or manual tests.
  
  \item \textbf{Stop}: \code{docker compose down} cleanly stops the service.
\end{enumerate}

This workflow is fast, repeatable, and suitable for CI/CD pipelines.

\subsection{Production Considerations}

For production deployments, additional considerations include:

\begin{enumerate}
  \item \textbf{Logging}: Aggregate containers logs to a centralized service (e.g., ELK, Splunk).
  
  \item \textbf{Monitoring}: Track metrics (CPU, memory, latency) using Prometheus/Grafana.
  
  \item \textbf{Orchestration}: Use Kubernetes for multi-node deployments and auto-scaling.
  
  \item \textbf{Secrets Management}: Use Docker secrets or external vaults (Vault, AWS Secrets Manager).
  
  \item \textbf{Image Registry}: Push images to a registry (Docker Hub, ECR, GCR) for distribution.
  
  \item \textbf{Network Policy}: Implement network policies to restrict inter-container and external communication.
  
  \item \textbf{Resource Limits}: Set CPU and memory limits to prevent resource exhaustion.
\end{enumerate}

These are deferred to operational governance.
